% Page budget: 1.55 pages | Owns: negation, OBC derivation, additional-state scope, implementation, and true-parameter condition.
% WRITING GUIDE
% Purpose: Explain negation, derive the final state-level orthogonal-by-construction map, and delimit its interpretability guarantee.
% Start here: Decisions D-185, D-190, D-192, D-194, D-198, D-202, D-203, D-222, D-229, and D-232 in docs/decisions.md.
% Supporting material: Current state-level OBC derivations under scripts/gantry/orthogonal-by-construction/documentation and their thesis-readiness guidance.
% Mandatory papers: Read Györök's projection-regularization paper and orthogonal-by-construction paper in full before writing; keep their guarantees and assumptions separate.
% Research-plan anchor: Preserve Aspect 3's goal of protecting physical interpretability; state clearly how the final construction differs from the proposed regularisation.
% Current decisions: Reduced identifiable coordinates, data-derived reference tuples, a tangent-only one-step physical-state projection, per-epoch refresh, and no affine comparison arm.
% Choices to justify: Protected parameter coordinates, reference-set construction, projected rows, rank tolerance, exclusion of the affine offset, refresh frequency, and treatment of additional states.
% Implementation audit: Read model_augmentation/fit_systems/obc.py, scripts/gantry/gantry_dynamic/obc_gantry.py, obc_diagnostics.py, and the OBC attachment in model.py before fixing the derivation.
% Reference check: Verify the original projection method and separate its claims from the state-level, nonlinear-parameter, LPV, and additional-state extensions developed here.
% Cautions: docs/orthogonal-projection-plan.md describes an older soft-penalty route; one-step state orthogonality is not trajectory-output orthogonality or guaranteed true recovery.
% Planned figures: New geometric projection figure with a physical and additional-state partition inset; show the remaining uniqueness limitation as clearly as the protected direction.
% Section is finished when: The protected space, coordinate frame, gradient treatment, added-state scope, and true-recovery condition are explicit.
\section{Interpretability by Orthogonal Construction}\label{sec:obc}

The split between $f_{\mathrm{base}}$ and $f_{\mathrm{aug}}$ in
\eqref{eq:aug_transition} is not unique (Section~\ref{sec:aug_structure}).
We remove this non-uniqueness from the one-step map by extending the
orthogonal-by-construction (OBC) parametrisation of Gy\"or\"ok et
al.~\cite{gyorok2026obc} from input-output models that are linear in their
parameters to the self-scheduled state transition of the gantry.
Section~\ref{sec:obc_negation} states the problem as a first-order
cancellation, Section~\ref{sec:obc_map} constructs the corrected transition,
and Section~\ref{sec:obc_scope} delimits its guarantee and states when it
recovers the true parameters.

\subsection{Negation of the Physical Parameters}\label{sec:obc_negation}
For a baseline that is nonlinear in its parameters, the changes that
$\theta_{\mathrm{base}}$ can produce are spanned, to first order, by the
columns of the transition sensitivity
\begin{equation}
 \Phi_{\bar\theta}(\tilde x,u)
 =\left.\frac{\partial f_{\mathrm{base}}(\theta_{\mathrm{base}},\tilde x,u)}
   {\partial\theta_{\mathrm{base}}}\right|_{\theta_{\mathrm{base}}=\bar\theta_{\mathrm{base}}}
   \in\R^{6\times10},
\label{eq:obc_sens}
\end{equation}
where $\bar\theta_{\mathrm{base}}$ is the expansion point, as
in~\cite[Eq.~(15)]{gyorok2025l4dc}.
To first order, a learned write that subtracts
$\Phi_{\bar\theta}\,\delta\theta$ compensates a parameter displacement
$\delta\theta$:
\begin{equation}
\begin{aligned}
 &f_{\mathrm{base}}(\bar\theta_{\mathrm{base}}+\delta\theta,\tilde x,u)
 +f_{\mathrm{aug}}(\theta_{\mathrm{aug}},\hat x,u)
 -\Phi_{\bar\theta}(\tilde x,u)\,\delta\theta\\
 &\quad\approx f_{\mathrm{base}}(\bar\theta_{\mathrm{base}},\tilde x,u)
 +f_{\mathrm{aug}}(\theta_{\mathrm{aug}},\hat x,u).
\end{aligned}
\label{eq:obc_negation}
\end{equation}
Both sides give the same one-step map, so the data cannot attribute such a
direction to either component.
Gy\"or\"ok et al.\ show the exact version for a baseline linear in its
parameters~\cite[Ex.~1]{gyorok2025l4dc}.
We call this cancellation negation.

On the gantry, the negated directions are generalised forces.
Differentiating \eqref{eq:eom} with respect to the $j$-th combination gives
\begin{equation}
 \frac{\partial\ddot q}{\partial\theta_{\mathrm{base},j}}
 =-M(Y)^{-1}\Bigl[\frac{\partial M(Y)}{\partial\theta_{\mathrm{base},j}}\,\ddot q
 +\frac{\partial C}{\partial\theta_{\mathrm{base},j}}\,\dot q
 +\frac{\partial K}{\partial\theta_{\mathrm{base},j}}\,q\Bigr],
\label{eq:obc_gantry_sens}
\end{equation}
which the RK4 step propagates into \eqref{eq:obc_sens}.
Combinations in $M(Y)$ act through $\ddot q$, those in $C$ through $\dot q$,
and $k_{b,\Sigma}$ through $q$.
A negating write therefore acts as an added inertia, damping or stiffness
force.
Through $M(Y)^{-1}$, its direction depends on the payload position.

OBC protects the ten combinations $\theta_{\mathrm{base}}$ of
Section~\ref{sec:params}.
The fourteen raw parameters would add the four null directions of that
section as dependent columns, while the construction below requires full
column rank~\cite[Assumption~1]{gyorok2026obc}.
Training uses the coordinates of Section~\ref{sec:params}, which map
invertibly to $\theta_{\mathrm{base}}$, so the sensitivity with respect to
them has the same range as \eqref{eq:obc_sens}.
The aim is a physical estimate that does not depend on the start of
training.
Whether it equals the true value is the subject of
Section~\ref{sec:obc_scope}.

\subsection{Orthogonal-by-Construction Transition}\label{sec:obc_map}
Gy\"or\"ok et al.\ make the network output orthogonal to the baseline
regressor by construction.
For an input-output baseline with stacked regressor $\Phi$, they subtract
from the stacked network output $F$ its least-squares fit on
$\Phi$~\cite[Eqs.~(8), (9), (11)]{gyorok2026obc}:
\begin{equation}
\begin{aligned}
 \theta_{\mathrm{aux}}&=(\Phi^\top\Phi)^{-1}\Phi^\top F,\qquad
 \hat Y=\Phi\theta_b+F-\Phi\,\theta_{\mathrm{aux}},\\
 0&=\Phi^\top\bigl(F-\Phi\,\theta_{\mathrm{aux}}\bigr),
\end{aligned}
\label{eq:obc_gyorok}
\end{equation}
where $\theta_b$ are the baseline parameters and $\hat Y$ the stacked
prediction.
The identity holds for every network parameter, so no weight is needed.
We therefore use this construction instead of the projection penalty
of~\cite[Eq.~(13)]{gyorok2025l4dc}, whose weight trades fit against
orthogonality~\cite[Sec.~1]{gyorok2026obc}.

We apply \eqref{eq:obc_gyorok} to the six physical rows of the normalised
transition, all of which $f_{\mathrm{aug}}$ writes.
Through $M(Y)^{-1}$ and the RK4 step, $f_{\mathrm{base}}$ has no
linear-in-parameters form.
The regressor is therefore the sensitivity \eqref{eq:obc_sens}, as in the
Taylor extension of~\cite[Sec.~4]{gyorok2025l4dc}.
Unlike~\cite[Eq.~(18)]{gyorok2025l4dc}, we do not append the affine offset
of that expansion.
The offset multiplies no parameter, and the parameter error depends only on
the part of the discrepancy in the range of the sensitivity
(Section~\ref{sec:obc_scope}).
Gy\"or\"ok et al.\ update the expansion point at every objective evaluation
or fix it at the nominal parameters~\cite[Sec.~4]{gyorok2025l4dc}.
Updating rebuilds the stack at every evaluation.
A fixed point would stay at the detuned start of training
(Section~\ref{sec:setup}).
We rebuild the stack at the current $\theta_{\mathrm{base}}$ at every epoch
boundary and hold it within the epoch.

The stack needs full states, but only positions are measured.
Gy\"or\"ok et al.\ name this as the open obstacle for state-space
baselines~\cite[Sec.~6]{gyorok2026obc}.
Their earlier fallback simulates the baseline on the training
inputs~\cite[Sec.~3]{gyorok2025l4dc}, but an open-loop simulation drifts on
the free axes $X$ and $Y$ (Section~\ref{sec:aug_closed_loop}).
We therefore build fixed reference tuples from every training sample with a
full stencil,
\begin{equation}
\begin{aligned}
 q_i&=P^{-\top}y_i,\qquad
 \dot q_i=\frac{-q_{i+2}+8q_{i+1}-8q_{i-1}+q_{i-2}}{12\,T_s},\\
 \tilde x_i&=\col(q_i,\dot q_i),\qquad \bar x_i=0,\qquad i=1,\dots,N,
\end{aligned}
\label{eq:obc_tuples}
\end{equation}
where $y_i$ are the recorded positions, paired with the recorded actuator
force $u_i$, and the stencil does not cross record boundaries.
Each tuple carries its measured $Y_i$, so the stack varies with the payload
position.
The recorded $u_i$ contains the feedback action.
The tuples are normalised as the training signals.
A fixed auxiliary evaluation set is allowed in~\cite[Sec.~3.2]{gyorok2026obc}.

On these tuples, \eqref{eq:obc_gyorok} becomes
\begin{equation}
\begin{aligned}
 \Phi_{\bar\theta}
 &=\col\bigl(\Phi_{\bar\theta}(\tilde x_1,u_1),\dots,
   \Phi_{\bar\theta}(\tilde x_N,u_N)\bigr)\in\R^{6N\times10},\\
 F_{\mathrm{aug}}
 &=\col\bigl(f_{\mathrm{aug}}(\theta_{\mathrm{aug}},\hat x_1,u_1),\dots,
   f_{\mathrm{aug}}(\theta_{\mathrm{aug}},\hat x_N,u_N)\bigr),\\
 \theta_{\mathrm{aux}}&=\Phi_{\bar\theta}^{+}F_{\mathrm{aug}},\\
 0&=\Phi_{\bar\theta}^\top
   \bigl(F_{\mathrm{aug}}-\Phi_{\bar\theta}\,\theta_{\mathrm{aux}}\bigr),\\
 0&=\Phi_{\bar\theta}^\top\,
   \frac{\partial\bigl(F_{\mathrm{aug}}-\Phi_{\bar\theta}\,
   \theta_{\mathrm{aux}}\bigr)}{\partial\theta_{\mathrm{aug}}},
\end{aligned}
\label{eq:obc_ours}
\end{equation}
where $\Phi_{\bar\theta}$ without arguments denotes the stack,
$\hat x_i=\col(\tilde x_i,\bar x_i)$, and $^{+}$ is the Moore-Penrose
pseudoinverse.
The fourth line is~\cite[Lemma~3]{gyorok2026obc}.
Since $\Phi_{\bar\theta}$ does not depend on $\theta_{\mathrm{aug}}$,
differentiating it gives the last line.
Every change of $\theta_{\mathrm{aug}}$ thus moves the corrected stack
orthogonally to the protected range.
This removes the flat direction of~\cite[Ex.~1]{gyorok2025l4dc} from the
one-step map.
$\theta_{\mathrm{aux}}$ is solved through a singular value decomposition of
the stack.
Its numerical rank is decided at the tolerance
$\max(6N,10)\,\epsilon\,\sigma_1$, where $\epsilon$ is the machine precision
and $\sigma_1$ the largest singular value.
A rebuild that loses full column rank keeps the previous stack.

We compute $\theta_{\mathrm{aux}}$ once per objective evaluation, from the
current network on the fixed tuples, and differentiate through this solve.
Fitting it on the rollout instead would make the correction at step $k$
depend on later rollout states.
The stack and $\bar\theta_{\mathrm{base}}$ are constants in the gradient.
The corrected augmented model is
\begin{equation}
\begin{aligned}
 \begin{bmatrix}\tilde x_{k+1}\\ \bar x_{k+1}\end{bmatrix}
 &=\begin{bmatrix}f_{\mathrm{base}}(\theta_{\mathrm{base}},\tilde x_k,\hat u_k)\\ 0\end{bmatrix}\\
 &\quad+\begin{bmatrix}f_{\mathrm{aug}}(\theta_{\mathrm{aug}},\hat x_k,\hat u_k)
   -\Phi_{\bar\theta}(\tilde x_k,\hat u_k)\,\theta_{\mathrm{aux}}\\
   g_{\mathrm{aug}}(\theta_{\mathrm{aug}},\hat x_k,\hat u_k)\end{bmatrix},
\end{aligned}
\label{eq:obc_model}
\end{equation}
which replaces the state equations of \eqref{eq:aug_transition}.
At every rollout point, the correction uses the expansion point of the
current epoch and the coefficient of the current objective, as Gy\"or\"ok
et al.\ predict with a fixed $\theta_{\mathrm{aux}}$~\cite[Eq.~(13)]{gyorok2026obc}.
Before each validation, $\theta_{\mathrm{aux}}$ is solved again for the
current network.
For the selected model, the stack is rebuilt at its $\theta_{\mathrm{base}}$
before this solve.
The correction $\Phi_{\bar\theta}(\tilde x_k,\hat u_k)\,\theta_{\mathrm{aux}}$
is a directional derivative of the RK4 step, computed by a forward
sensitivity through its stages (appendix).
The baseline has no added-state rows, so $g_{\mathrm{aug}}$ is not
corrected.

\subsection{Scope and Recovery Condition}\label{sec:obc_scope}
The orthogonality in \eqref{eq:obc_ours} holds for the one-step stack over
the tuples, at the expansion point.
It is a property of the sum over the tuples, not of each tuple.
Its inner product is Euclidean in normalised state coordinates, so the
state normalisation defines its metric.
The tuples lie on $\bar x=0$, so the coefficient does not see how
$f_{\mathrm{aug}}$ uses the added states.
The construction therefore does not constrain that use.
Off this slice, $f_{\mathrm{aug}}$ depends on added states that
$g_{\mathrm{aug}}$ set one step earlier.
Hence the simulated output trajectories are not made orthogonal.
The consistency and zero-covariance results
of~\cite[Thms.~16, 18]{gyorok2026obc} assume an objective that is the
prediction error of the protected stack itself.
Our objective \eqref{eq:aug_loss} is a closed-loop multistep output error,
so these results do not transfer.

A unique split need not be the true split.
Let $\theta^*_{\mathrm{base}}$ be the true parameters and $\Delta^*$ the
true one-step discrepancy on the tuples,
\begin{equation}
 \Delta^*=\col\bigl(\tilde x_{i+1}-f_{\mathrm{base}}(\theta^*_{\mathrm{base}},
 \tilde x_i,u_i)\bigr)_{i=1}^{N},
\label{eq:obc_delta}
\end{equation}
where the recorded successor $\tilde x_{i+1}$ is built with the stencil of
\eqref{eq:obc_tuples} and both terms are normalised.
We assume noise-free data and a corrected model that reproduces every
recorded successor, the counterpart of~\cite[Assumption~6]{gyorok2026obc}.
We further assume that $f_{\mathrm{base}}$ is affine in
$\theta_{\mathrm{base}}$ between the estimate $\hat\theta_{\mathrm{base}}$
and $\theta^*_{\mathrm{base}}$, and denote by $\Phi$ the stack at
$\theta^*_{\mathrm{base}}$.
The one-step equality on the tuples then reads
$\Phi(\hat\theta_{\mathrm{base}}-\theta^*_{\mathrm{base}})+\tilde F_{\mathrm{aug}}=\Delta^*$,
with $\tilde F_{\mathrm{aug}}=F_{\mathrm{aug}}-\Phi\,\theta_{\mathrm{aux}}$.
Multiplying by $\Phi^\top$ removes $\tilde F_{\mathrm{aug}}$ and gives
\begin{equation}
 \hat\theta_{\mathrm{base}}-\theta^*_{\mathrm{base}}=\Phi^{+}\Delta^*,
 \qquad
 \hat\theta_{\mathrm{base}}=\theta^*_{\mathrm{base}}
 \;\Longleftrightarrow\;
 \Phi^\top\Delta^*=0,
\label{eq:obc_bias}
\end{equation}
the state-level counterpart of~\cite[Cond.~4, Eq.~(21)]{gyorok2026obc}.
The model is trained on \eqref{eq:aug_loss}, so \eqref{eq:obc_bias} is a
one-step surrogate for the bias of the trained model, not that bias.

The condition is measured by the share of a stacked vector $v$ in the
protected range,
\begin{equation}
 \rho(v)=\frac{\norm{\Phi\Phi^{+}v}}{\norm{v}},
\label{eq:obc_overlap}
\end{equation}
where $\Phi$ is the stack of \eqref{eq:obc_bias}.
$\rho(\Delta^*)=0$ exactly when the condition holds.
Since $\Phi^{+}=\Phi^{+}\Phi\Phi^{+}$, the error in \eqref{eq:obc_bias} is
at most $\rho(\Delta^*)\norm{\Delta^*}/\sigma_{\min}(\Phi)$.
$\Delta^*$ requires $\theta^*_{\mathrm{base}}$, so $\rho(\Delta^*)$ is
available in simulation only.
Whether the condition holds depends on the excitation as well as on the
missing dynamics~\cite[Sec.~4.1]{gyorok2026obc}.
An input can be designed to satisfy it when the structure of the
discrepancy is known~\cite[Remark~5]{gyorok2026obc}.
Section~\ref{sec:results} therefore tests whether the estimate is
independent of the start.
It tests recovery of the true values only where the condition holds.
